% ============================================================================
% CAPITOLO 3 — PROGETTAZIONE DEL SISTEMA
% ============================================================================
\chapter{Progettazione del sistema}
\label{cap:sistema}

Questo capitolo illustra la fase di analisi e progettazione architetturale del sistema proposto. Lo studente deve mostrare la capacità di trasformare i requisiti teorici in un'architettura software strutturata, modulare e formalmente definita.

Inoltre, questo capitolo fornisce gli esempi pratici per l'inserimento e il corretto riferimento incrociato di:
\begin{itemize}
    \item \textbf{Figure e diagrammi} (\Cref{sec:esempi-figure});
    \item \textbf{Tabelle tipografiche con il pacchetto \texttt{booktabs}} (\Cref{sec:esempi-tabelle});
    \item \textbf{Formule ed equazioni matematiche numerate} (\Cref{sec:esempi-equazioni}).
\end{itemize}


% ============================================================================
\section{Analisi dei Requisiti}
\label{sec:requisiti}
% ============================================================================

\textbf{Guida per lo studente:} Prima di presentare l'architettura, descrivere i requisiti del sistema separandoli chiaramente in:
\begin{enumerate}
    \item \textbf{Requisiti Funzionali (RF):} Quali operazioni e funzionalità specifiche il sistema deve eseguire (es. RF1: Rilevamento di valori anomali sui sensori; RF2: Calcolo del punteggio di confidenza in tempo reale; RF3: Generazione automatica dei log di audit).
    \item \textbf{Requisiti Non Funzionali (RNF):} Proprietà qualitative del sistema, quali:
    \begin{itemize}
        \item \emph{Prestazioni e latenza:} Tempo massimo consentito per elaborare un singolo campione di input (es. $< 100$\,ms);
        \item \emph{Trasparenza e tracciabilità (White-box):} Esposizione di tutte le variabili intermedie per consentire l'ispezione;
        \item \emph{Modularità ed estensibilità:} Facilità di sostituzione dei moduli interni o aggiunta di nuovi scenari di test;
        \item \emph{Riproducibilità:} Garanzia che a parità di seme casuale (\emph{random seed}) e configurazione, i risultati siano identici al 100\%.
    \end{itemize}
\end{enumerate}


% ============================================================================
\section{Architettura Generale del Sistema}
\label{sec:architettura-generale}
% ============================================================================

\textbf{Guida per lo studente:} Descrivere l'architettura complessiva ad alto livello, spiegando come i vari macro-componenti cooperano tra loro e qual è il flusso dei dati.


% ----------------------------------------------------------------------------
\subsection{Esempio di inserimento di figure in LaTeX}
\label{sec:esempi-figure}
% ----------------------------------------------------------------------------

Ogni figura inserita nella tesi deve:
\begin{enumerate}
    \item Essere racchiusa nell'ambiente \texttt{figure} con posizionatore (es. \texttt{[htbp]});
    \item Avere un file sorgente in formato vettoriale (PDF) o raster ad alta risoluzione (PNG, JPEG da almeno 300 DPI) memorizzato nella cartella \texttt{img/};
    \item Essere provvista di una didascalia esplicativa (\texttt{\textbackslash caption\{\dots\}});
    \item Avere un'etichetta univoca (\texttt{\textbackslash label\{fig:\dots\}}) posizionata \textbf{sempre dopo} il caption;
    \item \textbf{Essere sempre citata nel testo} con \texttt{\textbackslash Cref\{fig:\dots\}} o \texttt{\textbackslash ref\{fig:\dots\}}. Una figura mai citata nel testo è un grave errore metodologico.
\end{enumerate}

La \Cref{fig:architettura-esempio} illustra a titolo di esempio uno schema architetturale e di pipeline (sostituibile con il diagramma UML o architetturale specifico della tesi), evidenziando le fasi di acquisizione, pre-elaborazione, estrazione delle feature e classificazione/decisione.

\begin{figure}[htbp]
    \centering
    \includegraphics[width=0.92\textwidth]{placeholder_figura.png}
    \caption{Esempio di schema architetturale e pipeline di elaborazione a blocchi (ispirato allo standard di riconoscimento e classificazione, Wikimedia Commons). Nel lavoro di tesi, sostituire questa figura con l'architettura specifica del proprio sistema o software.}
    \label{fig:architettura-esempio}
\end{figure}

% ----------------------------------------------------------------------------
\subsection{Esempio di diagramma di flusso (Flowchart)}
\label{sec:esempi-flowchart}
% ----------------------------------------------------------------------------

Accanto ai diagrammi architetturali (che mostrano la struttura statica dei moduli o delle pipeline), i \textbf{diagrammi di flusso (flowchart)} sono fondamentali per rappresentare visivamente la dinamica temporale, le diramazioni condizionali (\texttt{if-else}) e la sequenza delle operazioni decisionali. 

Un diagramma di flusso chiaro ed elegante \textbf{sostituisce con enorme efficacia lunghe pagine di codice sorgente}: permette al docente e alla commissione di cogliere la logica algoritmica in pochi istanti, senza perdersi nei dettagli implementativi di basso livello.

La \Cref{fig:diagramma-flusso} riporta un classico diagramma di flusso a norma ISO, illustrando l'uso standard dei blocchi di inizio/fine (ovali), di controllo condizionale (rombi di decisione con rami \textit{Yes/No}) e di elaborazione/azione (rettangoli).

\begin{figure}[htbp]
    \centering
    \includegraphics[width=0.45\textwidth]{placeholder_flusso.png}
    \caption{Esempio canonico di diagramma di flusso decisionale conforme agli standard grafici di processo (Wikimedia Commons). Nella tesi, sostituire con il flusso decisionale o computazionale del proprio algoritmo.}
    \label{fig:diagramma-flusso}
\end{figure}


% ============================================================================
\section{Componenti e Moduli Software}
\label{sec:moduli-sistema}
% ============================================================================

\textbf{Guida per lo studente:} Descrivere i singoli moduli software del sistema, le loro responsabilità e le interfacce esposte. Nel contesto del progetto di tesi, i componenti principali sono:

\begin{itemize}
    \item \textbf{Modulo Ambiente e Sensori (\texttt{Environment}):} Modella la dinamica fisica e lo strato sensoriale, introducendo rumore gaussiano o simulando iniezioni di dati malevoli;
    \item \textbf{Modello Generativo Interno (\texttt{GenerativeModel}):} Codifica la conoscenza a priori del sistema e genera predizioni sullo stato atteso;
    \item \textbf{Modulo di Stima dello Stato (\texttt{BeliefState}):} Calcola l'errore di predizione tra valore osservato e valore previsto, aggiornando la stima di incertezza posteriore;
    \item \textbf{Controllore Decisionale (\texttt{Controller}):} Valuta lo spazio delle azioni possibili attraverso una funzione di costo formale ed emette il comando attuativo.
\end{itemize}


% ----------------------------------------------------------------------------
\subsection{Linee guida ed esempio per la composizione di tabelle scientifiche}
\label{sec:esempi-tabelle}
% ----------------------------------------------------------------------------

Nelle pubblicazioni scientifiche e nelle tesi di laurea in Informatica, le tabelle devono rispettare precisi canoni tipografici:
\begin{enumerate}
    \item \textbf{Nessuna linea verticale:} Non usare mai linee verticali (\texttt{|}); rendono la lettura pesante e disordinata;
    \item \textbf{Pacchetto \texttt{booktabs}:} Usare esclusivamente le tre linee orizzontali ad alto valore tipografico: \texttt{\textbackslash toprule} (inizio), \texttt{\textbackslash midrule} (separatore intestazione) e \texttt{\textbackslash bottomrule} (chiusura);
    \item \textbf{Allineamento:} Allineare a sinistra (\texttt{l}) i testi descrittivi, al centro (\texttt{c}) i codici o stati, e a destra (\texttt{r}) i valori numerici;
    \item \textbf{Unità di misura nell'intestazione:} Inserire sempre l'unità di misura (es. ms, km/h, \%) nell'intestazione della colonna, mai ripetuta in ogni singola cella;
    \item \textbf{Didascalia e posizionamento:} La didascalia (\texttt{\textbackslash caption}) va posizionata preferibilmente \textbf{in alto} per le tabelle, seguita subito dal comando \texttt{\textbackslash label\{tab:\dots\}};
    \item \textbf{Riferimento obbligatorio:} Ogni tabella deve essere citata nel testo con \texttt{\textbackslash Cref\{tab:\dots\}}.
\end{enumerate}

La \Cref{tab:parametri-esempio} mostra una tabella tipografica standard compilata secondo queste direttive.

\begin{table}[htbp]
\centering
\small
\caption{Parametri di configurazione del sistema software.}
\label{tab:parametri-esempio}
\renewcommand{\arraystretch}{1.3}
\begin{tabular}{l c r l}
\toprule
\textbf{Parametro} & \textbf{Tipo} & \textbf{Default} & \textbf{Descrizione} \\
\midrule
\texttt{NUM\_STEPS}          & Intero  & 50    & Durata temporale della simulazione (passi) \\
\texttt{SENSOR\_NOISE}       & Float   & 0.05  & Ampiezza del rumore additivo uniforme ($\pm$) \\
\texttt{ANOMALY\_THRESHOLD}  & Float   & 0.30  & Soglia sul prediction error per anomalia \\
\texttt{DEFAULT\_SPEED}      & Float   & 10.0  & Velocità di marcia nominale (km/h) \\
\texttt{REDUCED\_SPEED}      & Float   & 4.0   & Velocità cautelativa con alta incertezza (km/h) \\
\texttt{STOP\_SPEED}         & Float   & 0.0   & Velocità di arresto immediato (km/h) \\
\bottomrule
\end{tabular}
\end{table}


% ============================================================================
\section{Modello Matematico e Logica Decisionale}
\label{sec:modello-matematico}
% ============================================================================

\textbf{Guida per lo studente:} Se la tesi impiega modelli probabilistici, equazioni differenziali, formule di costo o funzioni obiettivo, formalizzarli con precisione matematica.


% ----------------------------------------------------------------------------
\subsection{Esempio di inserimento di equazioni matematiche}
\label{sec:esempi-equazioni}
% ----------------------------------------------------------------------------

Le equazioni matematiche importanti devono essere posizionate nell'ambiente \texttt{equation}, dotate di etichetta \texttt{\textbackslash label\{eq:\dots\}} e referenziate con \texttt{\textbackslash Cref\{eq:\dots\}}.

Ad esempio, l'errore di predizione sensoriale $e_t$ al tempo $t$ è calcolato come discrepanza assoluta tra l'osservazione $o_t$ e la predizione $\hat{s}_t$:
\begin{equation}
e_t = \left| o_t - \hat{s}_t \right|
\label{eq:prediction-error}
\end{equation}

Quando l'errore di predizione $e_t$ definito nell'\Cref{eq:prediction-error} eccede la soglia $\theta_{\mathrm{anomalia}}$, il livello di incertezza $\sigma_t$ viene elevato al valore massimo.

La selezione della politica decisionale ottima $\pi^*$ avviene minimizzando la funzione obiettivo di Expected Free Energy (EFE), composta da una componente pragmatica e da una componente epistemica~\cite{friston2015epistemic}:
\begin{equation}
\pi^* = \arg\min_{\pi} \left[ -\mathrm{PragmaticValue}(\pi) - \mathrm{EpistemicValue}(\pi) \right]
\label{eq:efe-decision}
\end{equation}


% ----------------------------------------------------------------------------
\subsection{Esempio di inserimento di algoritmi in pseudocodice formale}
\label{sec:esempi-algoritmi}
% ----------------------------------------------------------------------------

Quando si descrive un procedimento computazionale in modo indipendente dallo specifico linguaggio di programmazione, è preferibile utilizzare lo **pseudocodice formale** tramite i pacchetti \texttt{algorithm} e \texttt{algpseudocode}.

I comandi principali dell'ambiente \texttt{algorithmic} sono:
\begin{itemize}
    \item \texttt{\textbackslash Require} e \texttt{\textbackslash Ensure}: Definiscono i dati in ingresso (input/precondizioni) e in uscita (output/postcondizioni);
    \item \texttt{\textbackslash State}: Introduce un'istruzione elementare o un'assegnazione ($\gets$ ottenuto con \texttt{\textbackslash gets});
    \item \texttt{\textbackslash If\{cond\} \dots \textbackslash Else \dots \textbackslash EndIf}: Blocchi condizionali;
    \item \texttt{\textbackslash For\{cond\} \dots \textbackslash EndFor} e \texttt{\textbackslash While\{cond\} \dots \textbackslash EndWhile}: Strutture iterative;
    \item \texttt{\textbackslash Call\{NomeFunzione\}\{Argomenti\}}: Invocazione di una procedura o metodo;
    \item \texttt{\textbackslash Comment\{testo\}}: Inserimento di commenti esplicativi a margine.
\end{itemize}

L'\Cref{alg:decision-process} illustra il ciclo di stima dello stato e decisione ottima eseguito a ogni passo temporale.

\begin{algorithm}[htbp]
\caption{Ciclo decisionale del controllore Active Inference.}
\label{alg:decision-process}
\begin{algorithmic}[1]
    \Require Osservazione sensoriale $o_t \in \mathbb{R}$, passo temporale $t \in \mathbb{N}$, soglia di anomalia $\theta \in \mathbb{R}^+$
    \Ensure Azione ottima selezionata $\pi^*$, credenza aggiornata $\hat{s}_t$
    \vspace{0.15cm}
    \State $\hat{s}_{\mathrm{pred}} \gets \Call{GenerativeModelPredict}{t}$ \Comment{Predizione dal prior interno}
    \State $e_t \gets |o_t - \hat{s}_{\mathrm{pred}}|$ \Comment{Calcolo del prediction error}
    \If{$e_t > \theta$}
        \State $\mathrm{anomalia} \gets \mathbf{True}$
        \State $\sigma_t \gets 1.0$ \Comment{Massima incertezza: sensore ritenuto inaffidabile}
        \State $\hat{s}_t \gets \hat{s}_{\mathrm{pred}}$ \Comment{L'agente si affida al modello generativo}
    \Else
        \State $\mathrm{anomalia} \gets \mathbf{False}$
        \State $\sigma_t \gets 0.05$ \Comment{Incertezza nominale (rumore di fondo)}
        \State $\hat{s}_t \gets o_t$ \Comment{L'agente acquisisce il dato dal sensore}
    \EndIf
    \For{\textbf{ciascuna} politica d'azione $\pi \in \Pi$}
        \State $\mathrm{val}_{\mathrm{pragmatic}} \gets \Call{ComputePragmaticValue}{\pi, \hat{s}_t}$
        \State $\mathrm{val}_{\mathrm{epistemic}} \gets \Call{ComputeEpistemicValue}{\pi, \sigma_t}$
        \State $\mathrm{EFE}(\pi) \gets -\mathrm{val}_{\mathrm{pragmatic}} - \mathrm{val}_{\mathrm{epistemic}}$
    \EndFor
    \State $\pi^* \gets \arg\min_{\pi \in \Pi} \mathrm{EFE}(\pi)$ \Comment{Selezione dell'azione che minimizza l'EFE}
    \State \Return $\pi^*, \hat{s}_t$
\end{algorithmic}
\end{algorithm}

Nel \Cref{cap:Implementazione} verrà esaminata l'implementazione concreta di tale algoritmo all'interno del codice Python sorgente.
